% 续接 ch02a：本文件从第 2 章 §2 末尾的引理 2.3 起，含 §3--§6。
上一段构造出的正合列其实是 $U^\circ$-模的正合列。显然，整个构造是函子式的，因此有下述结果。

\begin{lemma}\label{lem:02.2.2.3}
存在从 $\mathcal{M}(\D_\lambda)$ 到 $U^\circ$-模的短正合列所成范畴的共变函子，它把任意 $V\in \mathcal{M}(\D_\lambda)$ 映为正合列
\[
0\longrightarrow V\longrightarrow G\longrightarrow V(p_\alpha)\longrightarrow 0.
\]
\end{lemma}

\section{$\D_\lambda$-模上同调的消没定理\index{xiaomeidingli@消没定理|hyperpage}}
\label{sec:02.2.3}

本节讨论拟凝聚 $\D_\lambda$-模上同调的几个消没结果。

\begin{theorem}\label{thm:02.2.3.1}
设 $V$ 是旗簇 $X$ 上的拟凝聚 $\D_\lambda$-模。则当 $i>n(\lambda)$ 时，上同调群 $H^i(X,V)$ 消没。
\end{theorem}

特别地，当 $\lambda\in\h^*$ 反支配时，这包含了所有 $i>0$ 的 $H^i(X,V)$ 的消没，即有下述推论。

\begin{corollary}\label{cor:02.2.3.2}
设 $\lambda\in\h^*$ 反支配。则函子 $\Gamma$ 是从 $\mathrm{M}_{qc}(\D_\lambda)$ 到 $\mathcal{M}(U_\theta)$ 的正合函子。
\end{corollary}

先证 \ref{cor:02.2.3.2}，再用关于 $n(\lambda)$ 的归纳完成 \ref{thm:02.2.3.1} 的证明。

设 $G$ 是任意 $\calO_X$-模，$\mu\in P(\Sigma)$ 是支配权。记 $F$ 为以 $\mu$ 为最高权的有限维不可约 $\g$-模。在 §2 中，我们给 $\calO_X$-模 $\mathcal{F}=\calO_X\otimes_\C F$ 定义了滤过 $(\mathcal{F}_i;0\le i\le\dim F)$，其中各项是局部自由 $\calO_X$-子模，且 $\mathcal{F}_1=\calO(\mu)$。因此有 $\calO_X$-模的单射
\[
i_G:G=G(-\mu)\otimes_{\calO_X}\calO(\mu)=G(-\mu)\otimes_{\calO_X}\mathcal{F}_1
\longrightarrow G(-\mu)\otimes_{\calO_X}\mathcal{F} .
\]

设 $\lambda\in\h^*$ 反支配，$V$ 是拟凝聚 $\D_\lambda$-模，$\varphi:G\to V$ 是 $\calO_X$-模的态射。它诱导态射 $\varphi(-\mu):G(-\mu)\to V(-\mu)$，以及
\[
\varphi(-\mu)\otimes 1:G(-\mu)\otimes_{\calO_X}\mathcal{F}\longrightarrow V(-\mu)\otimes_{\calO_X}\mathcal{F} .
\]
此外还有自然的嵌入 $i_G:G\to G(-\mu)\otimes_{\calO_X}\mathcal{F}$ 与 $i_V:V\to V(-\mu)\otimes_{\calO_X}\mathcal{F}$，使图表
\begin{center}
\begin{tikzcd}[row sep=2.6em, column sep=5em]
G \arrow[r,"\varphi"] \arrow[d,"i_G"'] & V \arrow[d,"i_V"]\\
G(-\mu)\otimes_{\calO_X}\mathcal{F} \arrow[r,"\varphi(-\mu)\otimes 1"] & V(-\mu)\otimes_{\calO_X}\mathcal{F}
\end{tikzcd}
\end{center}
交换。于是在上同调层面有
\[
H^i(\varphi(-\mu)\otimes1)\circ H^i(i_G)=H^i(i_V)\circ H^i(\varphi),
\qquad 0\le i\le\dim X.
\]
又因 $\mathcal{F}$ 是自由 $\calO_X$-模，有
\[
H^i(X,G(-\mu)\otimes_{\calO_X}\mathcal{F})=H^i(X,G(-\mu))\otimes_\C F .
\]
再设 $G$ 是凝聚 $\calO_X$-模。可逆 $\calO_X$-模 $\calO(-2\rho)$ 是丰富\index{fengfu@丰富|hyperpage}的，因此可以取到支配权 $\mu$，使得 $1\le i\le\dim X$ 时 $H^i(X,G(-\mu))=0$。于是对这样的 $\mu\in P(\Sigma)$，有
\[
H^i(i_V)\circ H^i(\varphi)=0,\qquad 1\le i\le\dim X.
\]
由 \ref{lem:02.2.2.1}，$V(-\mu)\otimes_{\calO_X}\mathcal{F}$ 是 $V$ 与其 $\zz(\g)$-不变补的直和，也就是说 $i_V$ 有左逆 $j_V:V(-\mu)\otimes_{\calO_X}\mathcal{F}\to V$。因此
\[
H^i(\varphi)=H^i(j_V)\circ H^i(i_V)\circ H^i(\varphi)=0,\qquad 1\le i\le\dim X.
\]

任意拟凝聚 $\calO_X$-模都是其凝聚子模的直极限 [7, I.6.9.9]，而上同调与直极限交换 [8, III.2.9]。把上面的结论对各个嵌入 $G\subset V$ 取直极限，即知恒等态射 $j:V\to V$ 满足 $H^i(j)=0$（$1\le i\le\dim X$）。于是 $1\le i\le\dim X$ 时 $H^i(X,V)=0$，\ref{cor:02.2.3.2} 证毕。

下面证 \ref{thm:02.2.3.1}，这要用到 \ref{lem:02.2.2.3}。设 $\lambda\in\h^*$ 且 $n(\lambda)=k$。则有正合列
\[
0\longrightarrow V\longrightarrow G\longrightarrow V(p_\alpha)\longrightarrow 0,
\]
其中
\[
G=(V(-s_\alpha w^{-1}\sigma)\otimes_{\calO_X}\mathcal{F})[\lambda].
\]
如前所述，$n(\lambda+p_\alpha)<k$ 且 $n(\lambda-s_\alpha w^{-1}\sigma)<k$。于是由归纳假设，
\[
H^i(X,G)=H^i(X,V(-s_\alpha w^{-1}\sigma)\otimes_{\calO_X}\mathcal{F})[\lambda]
=\bigl(H^i(X,V(-s_\alpha w^{-1}\sigma))\otimes_\C F\bigr)[\lambda]=0
\]
以及
\[
H^i(X,V(p_\alpha))=0
\]
对 $i>k-1$ 成立。把上同调的长正合列用于上述短正合列，即得 $i>k$ 时 $H^i(X,V)=0$。

\section{$\D_\lambda$-模上同调的非消没定理\index{feixiaomeidingli@非消没定理|hyperpage}}
\label{sec:02.2.4}

设 $\lambda\in\h^*$，$\theta=W\cdot\lambda$。拟凝聚 $\D_\lambda$-模的范畴 $\mathrm{M}_{qc}(\D_\lambda)$ 是 $\mathcal{M}(\D_\lambda)$ 的厚子范畴。因此可以考虑 $\D_\lambda$-模的导出范畴\index{daochufanchou@导出范畴|hyperpage} $D(\mathcal{M}(\D_\lambda))$ 中的全子范畴 $D_{qc}(\mathcal{M}(\D_\lambda))$，它由「上同调为拟凝聚」的复形组成。设 $D(U_\theta)$ 是 $U_\theta$-模的导出范畴。范畴 $\mathcal{M}(\D_\lambda)$ 有足够多的内射对象，且内射对象都是 flasque 层。此外，整体截面函子 $\Gamma$ 的右上同调维数不超过 $\dim X$。因此可以定义导出函子\index{daochuhanzi@导出函子|hyperpage}
\[
R\Gamma:D(\mathcal{M}(\D_\lambda))\longrightarrow D(U_\theta).
\]

本节的主要目标是下述定理。

\begin{theorem}\label{thm:02.2.4.1}
设 $\lambda\in\h^*$ 正则，$C^\cdot,D^\cdot\in D_{qc}(\mathcal{M}(\D_\lambda))$，$f:C^\cdot\to D^\cdot$ 是态射。则下述条件等价：
\begin{enumerate}[label=(\roman*)]
  \item $f$ 是拟同构；
  \item $R\Gamma(f)$ 是拟同构。
\end{enumerate}
\end{theorem}

显然，(i) 蕴含 (ii)。反向的蕴涵来自 \ref{thm:02.2.4.1} 的下述特例。

\begin{lemma}\label{lem:02.2.4.2}
设 $\lambda\in\h^*$ 正则，$C^\cdot\in D_{qc}(\mathcal{M}(\D_\lambda))$ 且 $R\Gamma(C^\cdot)=0$。则 $C^\cdot=0$。
\end{lemma}

先说明 \ref{lem:02.2.4.2} 蕴含 \ref{thm:02.2.4.1}。设 $C^\cdot_f$ 是 $f$ 的映射锥，则有标准三角
\[
C^\cdot_f\longrightarrow C^\cdot\xrightarrow{\ f\ }D^\cdot\longrightarrow C^\cdot_f[1].
\]
对它施用 $R\Gamma$，得到特异三角
\[
R\Gamma(C^\cdot_f)\longrightarrow R\Gamma(C^\cdot)\xrightarrow{R\Gamma(f)}R\Gamma(D^\cdot)\longrightarrow R\Gamma(C^\cdot_f)[1].
\]
若 $R\Gamma(f)$ 是拟同构，则由上同调的长正合列可知 $i\in\Z$ 时 $H^i(R\Gamma(C^\cdot_f))=0$，即 $R\Gamma(C^\cdot_f)=0$。由 \ref{lem:02.2.4.2} 得 $C^\cdot_f=0$，故 $f$ 是拟同构。

下证 \ref{lem:02.2.4.2}。证明是对 $n(\lambda)$ 的归纳。

设 $n(\lambda)=0$。则由 \ref{cor:02.2.3.2}，$\Gamma$ 在 $\mathrm{M}_{qc}(\D_\lambda)$ 上正合。又不妨设 $C^\cdot$ 由 $\Gamma$-零调的 $\D_\lambda$-模组成。此时 $R\Gamma(C^\cdot)=\Gamma(C^\cdot)$。

设有某个 $i$ 使 $H^i(C^\cdot)\neq0$；它是拟凝聚 $\D_\lambda$-模。由于 $\calO(-2\rho)$ 丰富，存在支配权 $\mu\in P(\Sigma)$，使 $H^i(C^\cdot)(-\mu)$ 有非零的整体截面。由 \ref{lem:02.2.2.2}，若记 $F$ 为以 $-\mu$ 为最低权的有限维不可约 $\g$-模，$\mathcal{F}=\calO_X\otimes_\C F$，则
\[
\Gamma(X,H^i(C^\cdot))\otimes_\C F=\Gamma(X,H^i(C^\cdot)\otimes_{\calO_X}\mathcal{F})\neq 0 .
\]
于是 $\Gamma(X,H^i(C^\cdot))\neq 0$。
\begin{annote}
此处用 \ref{lem:02.2.2.2} 推出非零性，中间省了一步。由 \ref{lem:02.2.2.2}，$H^i(C^\cdot)(-\mu)$ 与 $H^i(C^\cdot)\otimes_{\calO_X}\mathcal{F}$ 只差一次 $\D$-模平移；平移不改变底层的 $\calO_X$-模，故后者也有非零的整体截面。又 $F$ 有限维，$\Gamma(X,H^i(C^\cdot)\otimes_{\calO_X}\mathcal{F})=\Gamma(X,H^i(C^\cdot))\otimes_\C F$，于是 $\Gamma(X,H^i(C^\cdot))\neq0$。[来源：底本 p.27，\ref{lem:02.2.2.2} 的陈述]
\end{annote}

另一方面，考虑短正合列
\[
0\longrightarrow\ker d^i\longrightarrow C^i\longrightarrow\im d^i\longrightarrow 0,
\]
与
\[
0\longrightarrow\im d^{i-1}\longrightarrow\ker d^i\longrightarrow H^i(C^\cdot)\longrightarrow 0.
\]
由上同调的长正合列可知
\[
H^n(X,\im d^i)=H^{n+1}(X,\ker d^i)\qquad (n\ge 1),
\]
以及
\[
H^n(X,\im d^{i-1})=H^n(X,\ker d^i)\qquad (n\ge 2).
\]
于是
\[
H^n(X,\im d^i)=H^{n+1}(X,\im d^{i-1})\qquad (n\ge 1).
\]
再由 $\Gamma$ 右上同调维数的有限性，$n>0$ 且 $i\in\Z$ 任意时 $H^n(X,\im d^i)=0$。这又给出 $n>0$、$i\in\Z$ 任意时 $H^n(X,\ker d^i)=0$。最终得到
\[
\Gamma(X,H^i(C^\cdot))=H^i(\Gamma(C^\cdot)),\qquad i\in\Z,
\]
这与 $H^i(\Gamma(C^\cdot))=0$ 矛盾。故所有 $i\in\Z$ 都有 $H^i(C^\cdot)=0$。

现设 $n(\lambda)=k>0$。同样不妨设 $C^\cdot$ 由 $\Gamma$-零调的 $\D_\lambda$-模组成，于是 $0=R\Gamma(C^\cdot)=\Gamma(C^\cdot)$。下面沿用 \ref{lem:02.2.2.3} 之前的记号。置 $\mu=\lambda-s_\alpha w^{-1}\sigma$。前文已证 $w\lambda$ 与 $w\mu$ 都反支配，且 $w\lambda-w\mu=s_{\beta_1}\sigma\in P(\Sigma)$。设 $\tau$ 是正则支配权，使 $\eta=\tau+w(\mu-\lambda)$ 支配。记 $F_{-\tau}$ 为以 $-\tau$ 为最低权的有限维表示，$\mathcal{F}_{-\tau}=\calO_X\otimes_\C F_{-\tau}$。则由 \ref{lem:02.2.2.2}，
\[
C^\cdot(-w^{-1}\tau)=(C^\cdot\otimes_{\calO_X}\mathcal{F}_{-\tau})[w\lambda-\tau],
\]
且
\[
H^i(X,C^j(-w^{-1}\tau))=H^i(X,C^j\otimes_{\calO_X}\mathcal{F}_{-\tau})[w\lambda-\tau]
=\bigl(H^i(X,C^j)\otimes_\C F_{-\tau}\bigr)[w\lambda-\tau]=0
\]
对 $i>0$ 成立，即复形 $C^\cdot(-w^{-1}\tau)$ 由 $\Gamma$-零调的 $\D_{\lambda-w^{-1}\tau}$-模组成。于是
\begin{align*}
R\Gamma(C^\cdot(-w^{-1}\tau))&=\Gamma(C^\cdot(-w^{-1}\tau))
=\Gamma(C^\cdot\otimes_{\calO_X}\mathcal{F}_{-\tau})[w\lambda-\tau]\\
&=\bigl(\Gamma(C^\cdot)\otimes_\C F_{-\tau}\bigr)[w\lambda-\tau]=0 .
\end{align*}
另一方面，取 $F^\eta$ 为以 $\eta$ 为最高权的有限维表示，$\mathcal{F}_\eta=\calO_X\otimes_\C F^\eta$。由 \ref{lem:02.2.2.1} 可得
\[
(C^\cdot(-w^{-1}\tau)\otimes_{\calO_X}\mathcal{F}_\eta)[\mu]=C^\cdot\bigl(-w^{-1}(\tau-\eta)\bigr)=C^\cdot(\mu-\lambda).
\]
重复前面的论证可知，$C^\cdot(\mu-\lambda)$ 由 $\Gamma$-零调的 $\D_\mu$-模组成，且
\[
\Gamma(C^\cdot(-s_\alpha w^{-1}\sigma))=\Gamma(C^\cdot(\mu-\lambda))=0 .
\]

现在令
\[
G^\cdot=(C^\cdot(-s_\alpha w^{-1}\sigma)\otimes_{\calO_X}\mathcal{F})[\lambda].
\]
显然 $G^\cdot$ 由 $\Gamma$-零调的 $U^\circ$-模组成，且 $\Gamma(G^\cdot)=0$。用 \ref{lem:02.2.2.3}，可知 $C^\cdot(p_\alpha)$ 由 $\Gamma$-零调的 $\D_{s_\alpha\lambda}$-模组成，且 $\Gamma(C^\cdot(p_\alpha))=0$。于是 $R\Gamma(C^\cdot(p_\alpha))=0$。按构造 $n(s_\alpha\lambda)<k$，故可用归纳假设，得 $C^\cdot(p_\alpha)=0$，最终 $C^\cdot=0$。\ref{lem:02.2.4.2} 证毕。

\begin{corollary}\label{cor:02.2.4.3}
设 $\lambda\in\h^*$ 正则，$V\in \mathrm{M}_{qc}(\D_\lambda)$，且其所有上同调模 $H^i(X,V)$（$i\in\Z_+$）都消没。则 $V=0$。
\end{corollary}

\begin{corollary}\label{cor:02.2.4.4}
设 $\lambda\in\h^*$ 反支配且正则。则任意 $V\in \mathrm{M}_{qc}(\D_\lambda)$ 都由其整体截面生成。
\end{corollary}

\begin{proof}
记 $W$ 为 $V$ 中由所有整体截面生成的 $\D_\lambda$-子模。则由 \ref{cor:02.2.3.2} 有 $U_\theta$-模的正合列
\[
0\longrightarrow\Gamma(X,W)\longrightarrow\Gamma(X,V)\longrightarrow\Gamma(X,V/W)\longrightarrow 0,
\]
故 $\Gamma(X,V/W)=0$。于是由 \ref{cor:02.2.4.3} 得 $V/W=0$，即 $V$ 由其整体截面生成。
\end{proof}

\section{Borel--Weil--Bott 定理}
\label{sec:02.2.5}

设 $n=\dim X$，$\lambda\in P(\Sigma)$，$\calO(\lambda)$ 是相应的可逆 $\calO_X$-模。则 $\calO(\lambda)$ 是 $\D_{\lambda-\rho}$-模，作为 $\calO_X$-模是凝聚的。由代数几何的一般结果（例如见 [8, Ch. III]）可知：
\begin{enumerate}[label=(\roman*)]
  \item $H^i(X,\calO(\lambda))$ 是有限维 $\g$-模；
  \item 若记 $\calO(\lambda)^{\vee}$ 为可逆层 $\calO(\lambda)$ 的对偶，$\omega_X$ 为 $X$ 上局部 $n$-形式构成的层，则由 Serre 对偶，向量空间 $H^i(X,\calO(\lambda))$ 的对偶同构于向量空间
\[
H^{n-i}(X,\calO(\lambda)^{\vee}\otimes_{\calO_X}\omega_X).
\]
\end{enumerate}
当然，$\calO(\lambda)^{\vee}\otimes_{\calO_X}\omega_X=\calO(-\lambda+2\rho)$。设 $w_0$ 是 $W$ 中的最长元。取 $w$ 为使 $w(\lambda-\rho)$ 反支配的最长元之一。则 $\ell(w)\ge n(\lambda-\rho)$，且等号成立当且仅当 $\lambda-\rho$ 不正则。另一方面，$w_0w(-\lambda+\rho)$ 也反支配，故由 [5, Ch. VI, Cor. 3 of Prop. 17]，
\[
n(-\lambda+\rho)\le\ell(w_0w)=\ell(w^{-1}w_0^{-1})=\ell(w^{-1}w_0)=\ell(w_0)-\ell(w)=n-\ell(w).
\]
于是
\begin{equation}\label{eq:02.2.5.1}
n(\lambda-\rho)\le\ell(w)\le n-n(-\lambda+\rho).
\end{equation}

若 $H^i(X,\calO(\lambda))\neq0$，则把 \ref{thm:02.2.3.1} 用于 $\calO(\lambda)$ 得 $i\le n(\lambda-\rho)$；若再用 (ii)，又得 $n-i\le n(-\lambda+\rho)$。因此
\begin{equation}\label{eq:02.2.5.2}
n-n(-\lambda+\rho)\le i\le n(\lambda-\rho).
\end{equation}
由 \eqref{eq:02.2.5.1} 与 \eqref{eq:02.2.5.2} 可见 $i=\ell(w)=n(\lambda-\rho)$。按前面的附注，这说明 $\lambda-\rho$ 正则。

下面只需在 $\lambda-\rho$ 正则时研究 $H^{\ell(w)}(X,\calO(\lambda))$。此时 $\mu=w(\lambda-\rho)$ 是正则反支配权。下文沿用 \ref{lem:02.2.2.3} 的记号与结果。置 $V=\calO(\lambda)$，则 $V(-s_\alpha w^{-1}\sigma)=\calO(\lambda-s_\alpha w^{-1}\sigma)$，而 $\lambda-s_\alpha w^{-1}\sigma-\rho$ 不正则。因此它的所有上同调群都消没。于是 $G$ 的所有上同调群也消没。所以正合列
\[
0\longrightarrow\calO(\lambda)\longrightarrow G\longrightarrow\calO(\lambda+p_\alpha)\longrightarrow 0
\]
给出：作为 $\g$-模，$i\in\Z_+$ 时 $H^i(X,\calO(\lambda+p_\alpha))=H^{i+1}(X,\calO(\lambda))$。现在 $p=-\check\alpha(\lambda-\rho)$，故
\[
\lambda+p\alpha=s_\alpha(\lambda-\rho)+\rho=w^{-1}s_{\beta_1}w(\lambda-\rho)+\rho=w^{-1}s_{\beta_1}\mu+\rho=w'^{-1}\mu+\rho .
\]
由关于 $\ell(w)$ 的归纳可得
\[
H^{\ell(w)}(X,\calO(\lambda))=\Gamma(X,\calO(\mu+\rho)).
\]
由 \ref{lem:02.2.2.2}，若 $F_{\mu+\rho}$ 是以 $\mu+\rho$ 为最低权的不可约 $\g$-模，且置 $\mathcal{F}_{\mu+\rho}=\calO_X\otimes_\C F_{\mu+\rho}$，则
\[
\calO(\mu+\rho)=(\mathcal{F}_{\mu+\rho})[\mu].
\]
于是
\begin{align*}
\Gamma(X,\calO(\mu+\rho))&=\Gamma(X,(\mathcal{F}_{\mu+\rho})[\mu])
=\Gamma(X,\calO_X\otimes_\C F_{\mu+\rho})[\mu]\\
&=\bigl(\Gamma(X,\calO_X)\otimes_\C F_{\mu+\rho}\bigr)[\mu]=F_{\mu+\rho}.
\end{align*}
至此 Borel--Weil--Bott 定理证毕。

\begin{theorem}\label{thm:02.2.5.1}
设 $\lambda\in P(\Sigma)$。则
\begin{enumerate}[label=(\roman*)]
  \item 若 $\lambda-\rho$ 不正则，则 $i\in\Z_+$ 时 $H^i(X,\calO(\lambda))=0$；
  \item 若 $\lambda-\rho$ 正则，且 $\mu\in P(\Sigma)$ 是反支配权，使对某个 $w\in W$ 有 $\lambda-\rho=w\mu$，则 $i\neq\ell(w)$ 时 $H^i(X,\calO(\lambda))=0$，而 $H^{\ell(w)}(X,\calO(\lambda))$ 是以 $\mu+\rho$ 为最低权的不可约有限维 $\g$-模。
\end{enumerate}
\end{theorem}

如 §1 那样，记 $N$ 为旗簇 $X$ 上由幂零根给出的向量丛，$\n^\circ$ 为 $N$ 的局部截面构成的局部自由 $\calO_X$-模。令
\[
W(j)=\{w\in W\mid\ell(w)=j\},\qquad 0\le j\le\dim X.
\]

\begin{lemma}\label{lem:02.2.5.2}
设 $0\le j\le\dim X$。则 $i\neq j$ 时 $H^i(X,\wedge^j\n^\circ)=0$，且 $H^j(X,\wedge^j\n^\circ)$ 是维数为 $\operatorname{Card}W(j)$ 的平凡 $\g$-模。
\end{lemma}

\begin{proof}
取基点 $x_0\in X$。显然，$\wedge^j\n^\circ$ 是 $G$-齐性向量丛 $\wedge^jN$ 的局部截面构成的 $\calO_X$-模，该丛由 Borel 子群 $B_{x_0}$ 在 $F_j=\wedge^j\n_{x_0}$ 上的自然表示确定。$B_{x_0}$ 在 $F_j$ 上的作用给出自然的 Jordan--H\"older 滤过：其各项是 $B_{x_0}$-不变子空间 $F_{jk}$（$0\le k\le\dim F_j$），满足 $\dim F_{jk}=k$，$\n_{x_0}F_{jk}\subset F_{jk-1}$，且 Cartan 代数 $\h$ 通过权 $\nu_{jk}$ 作用在 $F_{jk}/F_{jk-1}$ 上。这个权是 $\Sigma^+$ 中 $j$ 个不同根之和。该滤过诱导出向量丛 $\wedge^jN$ 的滤过，其中各项都是 $G$-齐性子丛。记 $\mathcal{F}_{jk}$（$0\le k\le\dim F_j$）为相应的局部截面构成的凝聚 $\calO_X$-模。显然 $\mathcal{F}_{jk}/\mathcal{F}_{jk-1}=\calO(\nu_{jk})$。

为计算 $\wedge^j\n^\circ$ 的上同调，需要更好地理解族 $\{\nu_{jk}\mid 0\le k\le\dim F_j,\ 0\le j\le\dim X\}$ 的结构。这个族等于 $\Sigma^+$ 的所有子集上的根之和。对每个 $\Phi\subset\Sigma^+$，记 $\nu(\Phi)$ 为 $\Phi$ 中所有根之和，并令
\[
S=\{\nu(\Phi)-\rho\mid\Phi\subset\Sigma^+\}.
\]
由于 $\nu(\Phi)-\rho$ 是 $\Phi$ 中根的半和与 $\Sigma^+-\Phi$ 中根的半和之差，显然 $S$ 在 Weyl 群 $W$ 的作用下不变。设 $S^-$ 是 $S$ 中反支配权的集合。则显然 $S=W\cdot S^-$。设 $\mu$ 是 $S^-$ 的正则元。记 $\omega_\alpha$ 为单根 $\alpha\in\Pi$ 对应的基本权。则 $\mu$ 是各 $\omega_\alpha$（$\alpha\in\Pi$）的线性组合，系数为严格负整数。因此 $\mu+\rho$ 仍反支配。另一方面，它又必须是正根之和。于是它等于 $0$，从而 $\mu=-\rho$。由此可见 $S$ 的正则元只有 $-w\rho$（$w\in W$）。此时
\[
\rho-w\rho=\nu\bigl(\Sigma^+\cap(-w(\Sigma^+))\bigr).
\]
又注意 $\ell(w)=\operatorname{Card}\bigl(\Sigma^+\cap(-w(\Sigma^+))\bigr)$。由 Borel--Weil--Bott 定理，若 $\nu_{jk}-\rho$ 奇异，则 $0\le i\le\dim X$ 时 $H^i(X,\calO(\nu_{jk}))=0$。反之，若 $\nu_{jk}-\rho$ 正则，则按前面的讨论 $\nu_{jk}=\rho-w\rho$，其中 $w\in W(j)$。此时
\[
H^i(X,\calO(\nu_{jk}))=0\quad(i\neq j),
\]
且
\[
H^j(X,\calO(\nu_{jk}))=H^j(X,\calO(\rho-w\rho))=\Gamma(X,\calO_X)=\C .
\]

利用这些信息与上同调的长正合列，对短正合列
\[
0\longrightarrow\mathcal{F}_{jk-1}\longrightarrow\mathcal{F}_{jk}\longrightarrow\calO(\nu_{jk})\longrightarrow 0
\]
作关于 $k$（$0\le k\le\dim F_j$）的归纳，即得所需结论。
\end{proof}

\section{$\D_\lambda$ 的上同调}
\label{sec:02.2.6}

本节要证明下述定理。

\begin{theorem}\label{thm:02.2.6.1}
\begin{enumerate}[label=(\roman*)]
  \item 自然映射 $U_\theta\to \D_\lambda$ 诱导 $U_\theta$ 到 $\Gamma(X,\D_\lambda)$ 的同构。
  \item $i>0$ 时 $H^i(X,\D_\lambda)=0$。
\end{enumerate}
\end{theorem}

设 $C^\cdot$ 是分次模 $U^\circ\otimes_{\calO_X}\wedge\n^\circ$，即对一切 $i\in\Z$ 有 $C^i=U^\circ\otimes_{\calO_X}\wedge^{-i}\n^\circ$。

首先注意，$C^\cdot$ 通过第一个因子上的左乘带有左 $\g$-模结构。外代数 $\wedge\n^\circ$ 有自然的左 $\g$-模结构。又 $U^\circ$ 通过右乘是右 $\g$-模，因此可以在 $C^\cdot$ 上定义另一个左 $\g$-模结构：
\[
\kappa(\xi)(u\otimes v)=-u\xi\otimes v+u\otimes\xi\cdot v,
\]
其中 $\xi\in\g$，$u\in U^\circ$，$v\in\wedge\n^\circ$。为说明这个定义有意义，考虑双加性映射 $\phi(\xi)$（$\xi\in\g$）。它从 $U^\circ\times\wedge\n^\circ$ 映入 $U^\circ\otimes_{\calO_X}\wedge\n^\circ$，由
\[
\phi(\xi)(u,v)=-u\xi\otimes v+u\otimes\xi\cdot v
\]
（$u\in U^\circ$，$v\in\wedge\n^\circ$）给出，则对任意 $f\in\calO_X$、$u\in U^\circ$ 与 $v\in\wedge\n^\circ$ 有
\begin{align*}
\phi(\xi)(uf,v)-\phi(\xi)(u,fv)
&=-(uf)\xi\otimes v+uf\otimes\xi\cdot v+u\xi\otimes fv-u\otimes\xi\cdot(fv)\\
&=u\xi(f)\otimes v+uf\otimes\xi\cdot v-u\otimes\xi(f)v-u\otimes f\xi\cdot v=0 .
\end{align*}
因此它经过 $U^\circ\otimes_{\calO_X}\wedge\n^\circ$ 分解，并诱导出 $\kappa(\xi)$。按构造，$C^\cdot$ 上的两个左 $\g$-模作用交换。因此可以把 $C^\cdot$ 看作左 $\g\times\g$-模：
\[
(\xi,\eta)w=\xi w+\kappa(\eta)w,\qquad \xi,\eta\in\g,\ w\in C^\cdot .
\]
此外，群 $G$ 作用在 $C^\cdot$ 上，其作用为 $U(\g)$ 上的伴随作用与 $\wedge\g$ 上的伴随作用的张量积。该作用的微分等于 $\g\times\g$-作用限制到对角子代数。因此可以把 $C^\cdot$ 看成 $(\g\times\g,G)$-模。考虑映射
\begin{align*}
d(u\otimes &v_1\wedge v_2\wedge\cdots\wedge v_k)\\
&=\sum_{i=1}^{k}(-1)^{i+1}uv_i\otimes v_1\wedge v_2\wedge\cdots\wedge\hat v_i\wedge\cdots\wedge v_k\\
&\quad+\sum_{1\le i<j\le k}(-1)^{i+j}u\otimes[v_i,v_j]\wedge v_1\wedge\cdots\wedge\hat v_i\wedge\cdots\wedge\hat v_j\\
&\qquad\wedge\cdots\wedge v_k,
\end{align*}
其中 $u\in U^\circ$，$v_1,v_2,\dots,v_k\in\n^\circ$。它是良定义的，因为 $\calO_X$ 的截面与 $\n^\circ$ 的截面在 $U^\circ$ 中交换；并且它把 $C^i$ 映入 $C^{i+1}$。又按定义，它与左乘给出的 $\g$-作用以及 $G$-作用交换。由于 $G$-作用的微分与左乘作用之差给出 $\kappa$ 所给出的作用，$d$ 是 $(\g\times\g,G)$-模的态射。直接计算还可验证 $d^2=0$，即 $C^\cdot$ 是 $(\g\times\g,G)$-模的复形。

\begin{lemma}\label{lem:02.2.6.2}
复形 $C^\cdot$ 零调。
\end{lemma}

\begin{proof}
先给复形 $C^\cdot$ 引入一个滤过。代数层 $U^\circ$ 有自然的滤过 $(F_pU^\circ;p\in\Z)$。令 $p<0$ 时 $F_pC^\cdot=0$，并对 $p\in\Z_+$ 置
\[
F_pC^\cdot=\sum_{q=-p}^{0}F_{p+q}U^\circ\otimes_{\calO_X}\wedge^{-q}\n^\circ .
\]
微分 $d$ 把 $F_pC^\cdot$ 映入自身（任意 $p\in\Z$），故 $C^\cdot$ 是滤过复形。相应的分次双复形为
\begin{align*}
\Gr_{p,q}C^\cdot&=F_p(U^\circ\otimes_{\calO_X}\wedge^{-q}\n^\circ)\big/F_{p-1}(U^\circ\otimes_{\calO_X}\wedge^{-q}\n^\circ)\\
&=S_{p+q}(\g^\circ)\otimes_{\calO_X}\wedge^{-q}\n^\circ,
\end{align*}
其微分 $\delta=\Gr d$ 是双次数 $(0,1)$ 的，由下式给出：
\[
\delta(u\otimes v_1\wedge v_2\wedge\cdots\wedge v_k)
=\sum_{i=1}^{k}(-1)^{i+1}uv_i\otimes v_1\wedge\cdots\wedge\hat v_i\wedge\cdots\wedge v_k,
\]
其中 $u\in S_{p-k}(\g^\circ)$，$v_1,v_2,\dots,v_k\in\n^\circ$。

取充分小的仿射开集 $U$，使 $\g^\circ|_U$ 与 $\n^\circ|_U$ 是自由 $\calO_U$-模。则由 Koszul 复形的标准结果 [3, Ch. X, §9, no. 3] 可知 $\Gr C^\cdot|_U$ 零调。因此 $\Gr C^\cdot$ 零调。

于是对任意 $p\in\Z$ 有复形的正合列
\[
0\longrightarrow F_{p-1}C^\cdot\longrightarrow F_pC^\cdot\longrightarrow\Gr_pC^\cdot\longrightarrow 0,
\]
而 $\Gr_pC^\cdot$ 零调。这说明对 $k\in-\N$ 与 $p\in\Z$ 有 $H^k(F_{p-1}C^\cdot)=H^k(F_pC^\cdot)$。又 $p<0$ 时 $F_pC^\cdot=0$，故
\[
H^k(F_pC^\cdot)=0,\qquad k\in-\N,\ p\in\Z,
\]
即所有 $F_pC^\cdot$ 都零调。

设 $\xi\in U^\circ\otimes_{\calO_X}\wedge^k\n^\circ$（$k>0$）满足 $d\xi=0$。由于 $C^\cdot$ 的滤过是穷尽的，存在 $p\in\Z_+$ 使 $\xi\in F_pC^\cdot$。于是由 $F_pC^\cdot$ 的零调性，存在 $\eta\in F_{p-k-1}U^\circ\otimes_{\calO_X}\wedge^{k+1}\n^\circ$ 使 $\xi=d\eta$。
\end{proof}

综合以上结果，得到下述命题。

\begin{proposition}\label{prop:02.2.6.3}
复形 $C^\cdot=U^\circ\otimes_{\calO_X}\wedge\n^\circ$ 是 $(\g\times\g,G)$-模 $\D_h$ 的一个左消解。
\end{proposition}

显然，对 $i,j\in\Z_+$，作为 $(\g\times\g,G)$-模有
\[
H^i(X,U^\circ\otimes_{\calO_X}\wedge^j\n^\circ)=H^i(X,U(\g)\otimes_\C\wedge^j\n^\circ)
=U(\g)\otimes_\C H^i(X,\wedge^j\n^\circ).
\]
由 \ref{lem:02.2.5.2}，$\g\times\g$ 在 $U(\g)\otimes_\C H^i(X,\wedge^j\n^\circ)$ 上的作用归结为 $U(\g)$ 的左右乘，即
\[
(\xi_1,\xi_2)\eta=\xi_1\eta-\eta\xi_2,\qquad \xi_1,\xi_2\in\g,\ \eta\in U(\g),
\]
群 $G$ 对 $U(\g)$ 的作用是伴随作用，对第二个因子的作用则是平凡的。此外，$i\neq j$ 时 $H^i(X,U^\circ\otimes_{\calO_X}\wedge^j\n^\circ)$ 消没，$i=j$ 时它是 $\operatorname{Card}W(j)$ 个 $U(\g)$ 的直和。于是计算 $\D_h$ 上同调的谱序列\index{puxulie@谱序列|hyperpage}（[6, 2.4, Remark 3]，它用上述消解）在 $E_2$ 项退化，从而
\begin{enumerate}[label=(\roman*)]
  \item $i>0$ 时 $H^i(X,\D_h)=0$；
  \item 视 $\Gamma(X,\D_h)$ 为 $(\g\times\g,G)$-模，它有有限递增滤过
\[
0=F_0\Gamma(X,\D_h)\subset F_1\Gamma(X,\D_h)\subset\cdots\subset F_n\Gamma(X,\D_h)=\Gamma(X,\D_h),
\]
使 $F_k\Gamma(X,\D_h)/F_{k-1}\Gamma(X,\D_h)$ 是 $\operatorname{Card}W(k)$ 个 $U(\g)$ 的直和，这些 $U(\g)$ 带上述 $(\g\times\g,G)$-模结构。
\end{enumerate}

按构造，滤过 $F\Gamma(X,\D_h)$ 是 $\Gamma(X,\D_h)$ 上自然 $G$-模结构的 $G$-模滤过，它在 $\Gamma(X,\D_h)$ 的 $G$-不变子代数上诱导滤过。由 \ref{lem:02.2.1.5}，这个子代数经 $\varphi$ 同构于 $U(\h)$。由于群 $G$ 半单且它在 $\Gamma(X,\D_h)$ 上的作用是代数的，$G$-模 $\Gamma(X,\D_h)$ 半单。于是 $\Gr\Gamma(X,\D_h)$ 的 $G$-不变元等于 $\Gr\varphi(U(\h))$。在上述 (ii) 中取 $G$-不变元，立得：
(iii) $F_k\varphi(U(\h))/F_{k-1}\varphi(U(\h))$ 是 $\operatorname{Card}W(k)$ 个 $\zz(\g)$ 的直和。

借助 Harish-Chandra 同态，可以把 $U(\h)$ 看作 $\zz(\g)$-模。这立得下述结果。

\begin{lemma}\label{lem:02.2.6.4}
泛包络代数 $U(\h)$ 是秩为 $\operatorname{Card}W$ 的自由 $\zz(\g)$-模。
\end{lemma}

另一方面，可以作张量积 $U(\g)\otimes_{\zz(\g)}U(\h)$，它带有自然的结合代数结构。它还有自然的 $G$-作用：在第一个因子上为伴随作用，在第二个因子上平凡。显然，$U(\h)$ 是该代数的 $G$-不变子代数。由 \ref{prop:02.2.1.6}，存在自然的代数同态
\[
\Psi:U(\g)\otimes_{\zz(\g)}U(\h)\longrightarrow\Gamma(X,\D_h),
\]
它由 $U(\g)$ 到 $\Gamma(X,\D_h)$ 的自然同态与 $\varphi$ 的张量积给出。

借助同构 $\varphi$，把 $\varphi(U(\h))$ 的滤过搬到 $U(\h)$ 上，并用
\[
F_p\bigl(U(\g)\otimes_{\zz(\g)}U(\h)\bigr)=U(\g)\otimes_{\zz(\g)}F_p(U(\h))
\]
定义 $U(\g)\otimes_{\zz(\g)}U(\h)$ 上的滤过。映射 $\Psi$ 显然与滤过相容。考虑相应的分次态射 $\Gr\Psi:U(\g)\otimes_{\zz(\g)}\Gr U(\h)\to\Gr\Gamma(X,\D_h)$。由前面的讨论，视 $\Gr\Gamma(X,\D_h)$ 为 $(\g\times\g,G)$-模，它是 $\operatorname{Card}W$ 个 $U(\g)$ 的直和。因此存在 $G$-不变元 $e_1,e_2,\dots,e_q$（$q=\operatorname{Card}W$），使
\[
\Gr\Gamma(X,\D_h)=\bigoplus_{1\le k\le q}U(\g)e_k,\qquad
\Gr\varphi(U(\h))=\bigoplus_{1\le k\le q}\zz(\g)e_k .
\]
于是 $\Gr\Psi$ 显然是同构，从而 $\Psi$ 也是同构。因此得到下述结果。

\begin{theorem}\label{thm:02.2.6.5}
\begin{enumerate}[label=(\roman*)]
  \item $\Gamma(X,\D_h)=U(\g)\otimes_{\zz(\g)}U(\h)$。
  \item $i>0$ 时 $H^i(X,\D_h)=0$。
\end{enumerate}
\end{theorem}

现在设 $V$ 是 $\h$-模。则由 \ref{prop:02.2.1.6} 有自然映射 $U(\g)\otimes_{\zz(\g)}V\to\Gamma(X,\D_h\otimes_{U(\h)}V)$。设
\[
\cdots\longrightarrow F^{-p}\longrightarrow F^{-p+1}\longrightarrow\cdots\longrightarrow F^{-1}\longrightarrow F^0\longrightarrow V\longrightarrow 0
\]
是 $V$ 的左自由 $U(\h)$-模消解。用 $\D_h$ 在 $U(\h)$ 上作张量积，得
\[
\cdots\longrightarrow \D_h\otimes_{U(\h)}F^{-p}\longrightarrow\cdots\longrightarrow \D_h\otimes_{U(\h)}F^0\longrightarrow \D_h\otimes_{U(\h)}V\longrightarrow 0 .
\]
由 \ref{cor:02.2.1.4}，$\D_h$ 局部 $U(\h)$-自由，故这是正合列。于是由 \ref{thm:02.2.6.5}(ii)，它是 $\D_h\otimes_{U(\h)}V$ 的、由 $\Gamma(X,-)$-零调层组成的左消解。这说明 $\D_h\otimes_{U(\h)}V$ 的所有高阶上同调消没。同时，利用 \ref{thm:02.2.6.5}(i)，得到正合列
\[
\cdots\longrightarrow U(\g)\otimes_{\zz(\g)}F^{-p}\longrightarrow\cdots\longrightarrow U(\g)\otimes_{\zz(\g)}F^0\longrightarrow\Gamma(X,\D_h\otimes_{U(\h)}V)\longrightarrow 0,
\]
结合 \ref{lem:02.2.6.4} 便得 $U(\g)\otimes_{\zz(\g)}V=\Gamma(X,\D_h\otimes_{U(\h)}V)$，且 $p\in\N$ 时 $\Tor_p^{\zz(\g)}(U(\g),V)=0$。因此有下述结果。

\begin{corollary}\label{cor:02.2.6.6}
设 $V$ 是任意 $U(\h)$-模。则
\begin{enumerate}[label=(\roman*)]
  \item $\Gamma(X,\D_h\otimes_{U(\h)}V)=U(\g)\otimes_{\zz(\g)}V$；
  \item $i>0$ 时 $H^i(X,\D_h\otimes_{U(\h)}V)=0$。
\end{enumerate}
\end{corollary}

特别地，若 $\lambda\in\h^*$，$\C_{\lambda+\rho}$ 是 $\h$ 通过 $\lambda+\rho$ 作用的一维 $\h$-模，则最终得到 \ref{thm:02.2.6.1}。
\begin{annote}
最后一步省略了向 $\D_\lambda$ 的过渡。由 §1，$\D_\lambda=\D_h/I_\lambda \D_h=\D_h\otimes_{U(\h)}\C_{\lambda+\rho}$，且 $U_\theta=U(\g)/J_\theta U(\g)=U(\g)\otimes_{\zz(\g)}\C_{\lambda+\rho}$。把 $V=\C_{\lambda+\rho}$ 代入 \ref{cor:02.2.6.6}，得 $\Gamma(X,\D_\lambda)=U(\g)\otimes_{\zz(\g)}\C_{\lambda+\rho}=U_\theta$，这就是 \ref{thm:02.2.6.1}(i)；上同调的消没就是 \ref{thm:02.2.6.1}(ii)。[来源：底本 p.20，§1 中 $\D_\lambda$ 与 $U_\theta$ 的定义]
\end{annote}
